\chapter*{Resumen}
\addcontentsline{toc}{chapter}{Resumen}

En esta tesis computacional se presenta el diseño geométrico cinemático, la simulación dinámica multipartícula y el análisis de estabilidad longitudinal para un acelerador lineal de partículas (\textit{LINAC}) de protones tipo Álvarez. La estructura resonante opera a una frecuencia de radiofrecuencia de $f_{\text{RF}} = 352.2\text{ MHz}$ y una fase síncrona nominal de $\phi_s = -30^\circ$, acelerando protones desde una energía inicial de $T_0 = 3.0\text{ MeV}$ hasta una energía de extracción de $T_f = 50.63\text{ MeV}$ a lo largo de $55$ celdas síncronas integradas en una cavidad de $10.59\text{ m}$.

Mediante un simulador modular desarrollado en Python, se caracterizó la atenuación del gradiente de campo eléctrico efectivo $E_{\text{acc}}$ (desde $11.24\text{ MV/m}$ en la celda 1 hasta $3.22\text{ MV/m}$ en la celda 55) y la tasa de ganancia por metro $dT/dz$. El seguimiento discreto del paquete de $N = 3000$ partículas en el espacio de fases $(\Delta\phi, \Delta T)$ demostró experimentalmente el Teorema de Liouville y el amortiguamiento adiabático de la dispersión de fase ($\sigma_\phi \propto (\beta\gamma)^{-3/4}$), la cual se comprimió de $15.2^\circ$ a $5.3^\circ$. Adicionalmente, el análisis de tolerancia delimitó una ventana operativa segura para errores de frecuencia RF de $\epsilon_f \le 1.0\%$ con capturas superiores al $98.2\%$, comparando dicho comportamiento con la dinámica ultrarelativista de electrones en celdas de longitud constante.

\vspace{0.5cm}
\noindent\textbf{Palabras clave:} Acelerador lineal, estructura Álvarez, sincronismo relativista, espacio de fases, amortiguamiento adiabático, Teorema de Liouville, estabilidad de fase.
